% Ch1 self-study -- "I do / You do" ladder for Zumdahl 1.1-1.10, four
% YOUR TURN questions per worked example. Values chosen disjoint from
% ch01-notes, ch01-ws1 (diagnostic) and ch01-ws2 (remediation).
%
% Two hazards paid for already in this chapter, per CLAUDE.md:
%  - imperial units are NOT siunitx units; write them as plain text
%    rather than declaring \mile / \inch for one document;
%  - every sig-fig item is checked against round-half ties, because a
%    worksheet teaching rounding must never sit ON the boundary.
\documentclass{apchem-worksheet}

\apcunitword{Chapter}
\apcunit{1}
\apcunittitle{Chemical Foundations}
\apcdoctitle{Self-Study \textbullet\ Chapter 1, I~do / You~do}
\apcsubtitle{Zumdahl \S1.1--1.10 \textbullet\ PDF pp.\ 31--57 \textbullet\ four YOUR TURN questions per ladder}

\begin{document}
\apctitleblock

\begin{apcnote}[How to use these notes --- read this first]
Each skill is a \textbf{ladder}: a fully worked example in the
solid-framed box, then \textbf{four} YOUR TURN questions in the dashed
box --- same skill, new numbers, no help.

\begin{enumerate}[leftmargin=*,itemsep=0.15em]
  \item \textbf{Work all four} before checking anything.
  \item Check against the gray \emph{check:} line. All four right on the
        first try earns the tick on the tracker (last page).
  \item Any misses: re-read the worked example, then redo the missed ones
        from a blank page.
\end{enumerate}

\textbf{Why this chapter deserves real effort.} Nothing here is
difficult, and that is exactly the trap: these are the habits every later
chapter assumes. A stoichiometry answer with the wrong number of
significant figures, or a gas law worked in \si{\celsius}, loses the mark
no matter how good the chemistry was. Get these seven skills automatic
now and they cost you nothing for the rest of the year.
\end{apcnote}

% =====================================================================
\section*{\sffamily Ladder 1 \textbullet\ Units, prefixes, and scientific
notation}
\zsec{1.3}

SI gives each quantity one base unit; prefixes scale it by powers of ten.

\begin{center}
\renewcommand{\arraystretch}{1.2}
\begin{tabular}{@{}llll@{}}
\hline
\textbf{Prefix} & \textbf{Symbol} & \textbf{Factor} &
  \textbf{Example} \\
\hline
kilo  & k & $10^{3}$    & \SI{1}{\kilo\gram} $= \SI{1000}{\gram}$ \\
centi & c & $10^{-2}$   & \SI{1}{\centi\meter} $= \SI{0.01}{\meter}$ \\
milli & m & $10^{-3}$   & \SI{1}{\milli\litre} $= \SI{0.001}{\litre}$ \\
micro & $\mu$ & $10^{-6}$ & \si{\micro\litre} \\
nano  & n & $10^{-9}$   & --- \\
\hline
\end{tabular}
\end{center}

Volume has a useful identity worth memorizing outright:
\[ \SI{1}{\milli\litre} = \SI{1}{\centi\meter\cubed} \]

\begin{workedexample}[I do: moving between prefixes]
\textbf{(a) Convert \SI{4.50}{\kilo\gram} to milligrams.} Go through the
base unit rather than guessing an exponent:
\[ 4.50~\text{kg} \times \frac{1000~\text{g}}{1~\text{kg}}
   \times \frac{1000~\text{mg}}{1~\text{g}}
   = \mathbf{4.50\times10^{6}~mg} \]

\textbf{(b) Convert \SI{250}{\nano\meter} to metres.}
\[ 250 \times 10^{-9} = \mathbf{2.50\times10^{-7}~m} \]

\textbf{(c) Express \SI{0.0000345}{\litre} in scientific notation and
in microlitres.} Move the point until one non-zero digit sits to its
left:
\[ 3.45\times10^{-5}~\text{L}
   = 3.45\times10^{-5} \times 10^{6}~\si{\micro\litre}
   = \mathbf{34.5}~\si{\micro\litre} \]

\textbf{The check that catches sign errors:} a \emph{smaller} unit must
give a \emph{bigger} number. Milligrams are tiny, so (a) had to grow;
metres are large, so (b) had to shrink.
\end{workedexample}

\begin{yourturn}[YOUR TURN 1 --- four questions]
\begin{enumerate}[leftmargin=*,itemsep=0.5em]
  \item \SI{3.20}{\centi\meter} in metres: \workspace[1.2cm]
  \item \SI{75}{\milli\litre} in litres, and in \si{\centi\meter\cubed}:
        \workspace[1.4cm]
  \item \SI{0.00845}{\gram} in scientific notation, and in milligrams:
        \workspace[1.4cm]
  \item \SI{2.6}{\kilo\meter} in centimetres: \workspace[1.4cm]
\end{enumerate}
\selfcheck{(a) \SI{0.0320}{\meter} \quad (b) \SI{0.075}{\litre} $=$
\SI{75}{\centi\meter\cubed} \quad (c) $8.45\times10^{-3}$~g $=$
\SI{8.45}{\milli\gram} \quad (d) $2.6\times10^{5}$~cm}
\keyonly{\begin{apcanswer}(a) $3.20 \times 10^{-2} =
\mathbf{0.0320}$~m. (b) $75 \times 10^{-3} = \mathbf{0.075}$~L; and
since \SI{1}{\milli\litre} $=$ \SI{1}{\centi\meter\cubed}, it is
\textbf{75}~\si{\centi\meter\cubed} with no arithmetic at all.
(c) $\mathbf{8.45\times10^{-3}}$~g $= \mathbf{8.45}$~mg.
(d) $2.6~\text{km} \times 1000 \times 100 =
\mathbf{2.6\times10^{5}}$~cm --- smaller unit, bigger
number.\end{apcanswer}}
\end{yourturn}

% =====================================================================
\section*{\sffamily Ladder 2 \textbullet\ Precision, accuracy, and
uncertainty}
\zsec{1.4}

Two words that are not synonyms:

\begin{center}
\renewcommand{\arraystretch}{1.2}
\begin{tabular}{@{}lll@{}}
\hline
 & \textbf{Definition} & \textbf{Poor when} \\
\hline
\textbf{Accuracy} & closeness to the \emph{true} value &
  there is a systematic error \\
\textbf{Precision} & closeness of repeats to \emph{each other} &
  there is random scatter \\
\hline
\end{tabular}
\end{center}

A miscalibrated balance gives readings that agree beautifully with each
other and are all wrong: \textbf{precise but inaccurate}. That
combination is the signature of \term{systematic error}, and no amount of
repeating fixes it --- only recalibration does.

Every measurement carries one \textbf{estimated} final digit. Reading a
burette marked every \SI{0.1}{\milli\litre}, you record
\SI{23.46}{\milli\litre}: three certain digits and one estimated.

\begin{workedexample}[I do: judging a data set]
Three students each mass the same object four times. The true mass is
\SI{4.000}{\gram}.
\begin{center}
\renewcommand{\arraystretch}{1.15}
\begin{tabular}{@{}llll@{}}
\hline
\textbf{Student} & \textbf{Readings (g)} & \textbf{Mean} &
  \textbf{Verdict} \\
\hline
A & 3.998, 4.001, 3.999, 4.002 & 4.000 & accurate \emph{and} precise \\
B & 4.512, 4.510, 4.511, 4.513 & 4.512 & precise, not accurate \\
C & 3.85, 4.15, 3.92, 4.08 & 4.00 & accurate on average, not precise \\
\hline
\end{tabular}
\end{center}

\textbf{Student B} is the instructive case: the readings agree to
\SI{0.001}{\gram}, so the technique is excellent, but every one is high
by about \SI{0.51}{\gram}. That is a systematic error --- a balance that
was not zeroed. Repeating the measurement a hundred times would not
reveal it; only checking against a known standard would.

\textbf{Student C} shows why a good mean is not proof of a good
experiment: the average lands on 4.00 by luck of cancellation, while no
individual reading is close.
\end{workedexample}

\begin{yourturn}[YOUR TURN 2 --- four questions]
\begin{enumerate}[leftmargin=*,itemsep=0.5em]
  \item A thermometer reads \SI{2.0}{\celsius} too high every time. Which
        is affected, accuracy or precision? \workspace[1.2cm]
  \item Readings 25.1, 25.9, 24.8, 25.6 with a true value of
        \SI{25.4}{\celsius}: describe both qualities.
        \workspace[1.4cm]
  \item A burette is marked every \SI{0.1}{\milli\litre}. How many
        decimal places should you record, and why? \workspace[1.4cm]
  \item Which type of error can be reduced by averaging more trials, and
        which cannot? \workspace[1.4cm]
\end{enumerate}
\selfcheck{(a) accuracy only \quad (b) accurate on average, imprecise
\quad (c) two --- estimate one digit past the smallest marking \quad
(d) random can be averaged down; systematic cannot}
\keyonly{\begin{apcanswer}(a) \textbf{Accuracy} --- the readings are
consistently offset, so they still agree with each other. This is
systematic error. (b) The mean is 25.35, close to 25.4, so it is
\textbf{accurate on average but imprecise} --- the spread is over
\SI{1}{\celsius}. (c) \textbf{Two} decimal places, e.g.
\SI{23.46}{\milli\litre}: read the marked digit, then estimate one
further digit. Every measurement ends in an estimated digit. (d)
\textbf{Random} error scatters both ways and averages toward zero;
\textbf{systematic} error shifts every reading the same way and survives
any number of repeats.\end{apcanswer}}
\end{yourturn}

% =====================================================================
\section*{\sffamily Ladder 3 \textbullet\ Counting significant figures}
\zsec{1.5}

\begin{enumerate}[leftmargin=*,itemsep=0.15em]
  \item Non-zero digits always count.
  \item Zeros \emph{between} non-zeros count (captive): 1005 has 4.
  \item Leading zeros never count --- they only place the decimal:
        0.0042 has 2.
  \item Trailing zeros count \textbf{only if a decimal point is
        present}: 100 has 1, but 100. has 3 and 100.0 has 4.
  \item \textbf{Exact numbers} --- counted objects and definitions --- have
        infinite significant figures and never limit an answer.
\end{enumerate}

\begin{workedexample}[I do: counting, including the awkward cases]
\begin{center}
\renewcommand{\arraystretch}{1.2}
\begin{tabular}{@{}lll@{}}
\hline
\textbf{Number} & \textbf{Sig figs} & \textbf{Why} \\
\hline
0.00560         & 3 & leading zeros place only; the trailing 0 counts \\
9004            & 4 & captive zeros count \\
2500            & 2 & no decimal point, so trailing zeros do not count \\
2500.           & 4 & the decimal point makes them count \\
$2.50\times10^{3}$ & 3 & scientific notation removes all ambiguity \\
12 eggs         & exact & counted, not measured \\
\SI{1}{\kilo\gram} $=$ \SI{1000}{\gram} & exact & a definition \\
\hline
\end{tabular}
\end{center}

\textbf{The lesson in rows 3--5:} ``2500'' is genuinely ambiguous in
ordinary writing. Scientific notation is not decoration --- it is how you
say precisely how many digits you mean.
\end{workedexample}

\begin{yourturn}[YOUR TURN 3 --- four questions]
Give the number of significant figures:
\begin{enumerate}[leftmargin=*,itemsep=0.5em]
  \item 0.03080 \quad and \quad 40.0 \workspace[1.2cm]
  \item 6000 \quad and \quad 6000. \workspace[1.2cm]
  \item $7.20\times10^{-4}$ \quad and \quad 1.0090 \workspace[1.2cm]
  \item How many significant figures does the ``60'' carry in
        ``\SI{1}{\minute} $=$ \SI{60}{\second}''? Explain.
        \workspace[1.4cm]
\end{enumerate}
\selfcheck{(a) 4 and 3 \quad (b) 1 and 4 \quad (c) 3 and 5 \quad
(d) exact --- infinite, because it is a definition}
\keyonly{\begin{apcanswer}(a) \textbf{0.03080} has 4 (leading zeros
place; captive and trailing zeros count); \textbf{40.0} has 3.
(b) \textbf{6000} has 1; \textbf{6000.} has 4.
(c) $7.20\times10^{-4}$ has 3 --- only the coefficient counts;
\textbf{1.0090} has 5. (d) \textbf{Exact}, so it never limits an answer:
the minute is \emph{defined} as sixty seconds, it was not
measured.\end{apcanswer}}
\end{yourturn}

% =====================================================================
\section*{\sffamily Ladder 4 \textbullet\ Calculating with significant
figures}
\zsec{1.5}

Two rules, and they are different. Mixing them is the single commonest
error in this chapter.

\begin{center}
\fbox{\parbox{0.92\linewidth}{\centering
\textbf{Multiply / divide} $\Rightarrow$ count \textbf{significant
figures}; the answer takes the fewest.\\[2pt]
\textbf{Add / subtract} $\Rightarrow$ count \textbf{decimal places}; the
answer takes the fewest.}}
\end{center}

And one procedural rule: in a multi-step calculation, \textbf{round only
at the end}. Rounding intermediates lets error accumulate.

\begin{workedexample}[I do: both rules, and a subtraction that surprises]
\textbf{(a) $6.221 \times 5.2$.} Multiplication, so count sig figs:
4 and 2, fewest is 2.
\[ 6.221 \times 5.2 = 32.3492 \to \mathbf{32} \]

\textbf{(b) $18.7 + 2.34 + 0.891$.} Addition, so count \emph{decimal
places}: 1, 2, 3 --- fewest is 1.
\[ 18.7 + 2.34 + 0.891 = 21.931 \to \mathbf{21.9} \]
Note the answer has 3 sig figs while one input had only 3 and another had
4 --- with addition, sig figs are not what governs.

\textbf{(c) $8.55 - 8.32$.} Both inputs have 3 sig figs, but subtraction
counts decimal places (2 each):
\[ 8.55 - 8.32 = \mathbf{0.23} \]
The answer has only \textbf{2} significant figures. Subtracting nearby
numbers destroys precision --- which is why a titration's
\emph{difference} of two burette readings is the weak link in the
measurement.
\end{workedexample}

\begin{yourturn}[YOUR TURN 4 --- four questions]
\begin{enumerate}[leftmargin=*,itemsep=0.5em]
  \item $4.12 \times 0.20$ \workspace[1.2cm]
  \item $105.6 + 3.27$ \workspace[1.2cm]
  \item $\dfrac{27.9}{1.3}$ \workspace[1.2cm]
  \item $12.00 - 11.61$ --- give the answer and say how many significant
        figures survived. \workspace[1.4cm]
\end{enumerate}
\selfcheck{(a) 0.82 \quad (b) 108.9 \quad (c) 21 \quad (d) 0.39, only
2 sig figs}
\keyonly{\begin{apcanswer}(a) $0.824 \to \mathbf{0.82}$ (2 sig figs,
from 0.20). (b) $108.87 \to \mathbf{108.9}$ (1 decimal place, from
105.6). (c) $21.46 \to \mathbf{21}$ (2 sig figs, from 1.3).
(d) $\mathbf{0.39}$ --- both inputs had 4 sig figs, but subtraction keeps
2 decimal places and only \textbf{2} significant figures survive. This is
loss of significance, and it is why you never subtract two large nearly
equal measurements if you can avoid it.\end{apcanswer}}
\end{yourturn}

% =====================================================================
\section*{\sffamily Ladder 5 \textbullet\ Dimensional analysis}
\zsec{1.6--1.7}

Write every conversion as a fraction, arranged so the unit you have
cancels. If the units come out right, the arithmetic is almost certainly
right too --- and if they do not, no amount of correct arithmetic will
save the answer.

\begin{workedexample}[I do: a two-step conversion, units first]
A car's fuel use is \SI{7.60}{\litre} per \SI{100}{\kilo\meter}. How many
litres for a \SI{450}{\kilo\meter} trip? Then: how many millilitres per
kilometre?

\textbf{Set up so kilometres cancel:}
\[ 450~\text{km} \times \frac{7.60~\text{L}}{100~\text{km}}
   = \mathbf{34.2~L} \]
The kilometres cancel top and bottom, leaving litres --- which is what
was asked, so the setup is right.

\textbf{Second part, chaining two factors:}
\[ \frac{7.60~\text{L}}{100~\text{km}}
   \times \frac{1000~\text{mL}}{1~\text{L}}
   = \mathbf{76.0~mL/km} \]

\textbf{Sense check:} 450 km is 4.5 times the 100 km reference, and
$4.5 \times 7.60 = 34.2$. The dimensional setup and the mental estimate
agree.
\end{workedexample}

\begin{yourturn}[YOUR TURN 5 --- four questions]
Use 1 inch $=$ \SI{2.54}{\centi\meter} (exact) where needed.
\begin{enumerate}[leftmargin=*,itemsep=0.5em]
  \item Convert 12.0 inches to centimetres: \workspace[1.2cm]
  \item Convert \SI{5.00}{\meter} to inches: \workspace[1.4cm]
  \item A tap fills \SI{2.5}{\litre} in \SI{40}{\second}. How long to
        fill \SI{15}{\litre}? \workspace[1.6cm]
  \item In question (b), which unit had to cancel, and where did you
        place it to make that happen? \workspace[1.4cm]
\end{enumerate}
\selfcheck{(a) \SI{30.5}{\centi\meter} \quad (b) 197 in \quad
(c) \SI{240}{\second} (4.0 min) \quad (d) centimetres, placed in the
denominator}
\keyonly{\begin{apcanswer}(a) $12.0 \times 2.54 = \mathbf{30.5}$~cm
(3 sig figs from 12.0; the 2.54 is exact).
(b) $5.00~\text{m} \times \frac{100~\text{cm}}{1~\text{m}} \times
\frac{1~\text{in}}{2.54~\text{cm}} = \mathbf{197}$~in.
(c) $15~\text{L} \times \frac{40~\text{s}}{2.5~\text{L}} =
\mathbf{240}$~s $= 4.0$~min.
(d) \textbf{Centimetres.} After converting metres to centimetres, cm sat
in the numerator, so the inch factor had to be written with cm in the
\textbf{denominator} to cancel it.\end{apcanswer}}
\end{yourturn}

% =====================================================================
\section*{\sffamily Ladder 6 \textbullet\ Temperature and density}
\zsec{1.8--1.9}

\[ T_{\text{K}} = T_{^\circ\text{C}} + 273.15 \qquad
   T_{^\circ\text{C}} = \frac{5}{9}\left(T_{^\circ\text{F}} - 32\right)
   \qquad d = \frac{m}{V} \]

Celsius and Kelvin degrees are the \emph{same size} --- only the zero
moves --- so a temperature \emph{difference} is numerically identical in
both. Fahrenheit degrees are a different size, which is why its
conversion needs both a factor and an offset.

Density is a conversion factor between mass and volume, usable in either
direction. Water is \SI{1.00}{\gram\per\milli\litre}, which makes it the
reference for whether something floats.

\begin{workedexample}[I do: temperature, then density both ways]
\textbf{(a) Convert \SI{-40}{\celsius} to kelvin and to Fahrenheit.}
\[ T_{\text{K}} = -40 + 273.15 = \mathbf{233~K} \]
\[ T_{^\circ\text{F}} = \frac{9}{5}(-40) + 32 = -72 + 32 =
   \mathbf{-40~^\circ F} \]
The famous coincidence: $-40$ is the one temperature where the two
scales read alike.

\textbf{(b) A metal cube of side \SI{2.00}{\centi\meter} has a mass of
\SI{86.4}{\gram}. Find its density.}
\[ V = (2.00)^3 = \SI{8.00}{\centi\meter\cubed} \qquad
   d = \frac{86.4}{8.00} = \mathbf{10.8~g/cm^3} \]

\textbf{(c) What volume would \SI{250.}{\gram} of that metal occupy?}
Run the density backward:
\[ V = \frac{m}{d} = \frac{250.}{10.8} = \mathbf{23.1~cm^3} \]
\end{workedexample}

\begin{yourturn}[YOUR TURN 6 --- four questions]
\begin{enumerate}[leftmargin=*,itemsep=0.5em]
  \item Convert \SI{25}{\celsius} to kelvin, and \SI{350}{\kelvin} to
        \si{\celsius}: \workspace[1.4cm]
  \item A liquid has mass \SI{45.0}{\gram} and volume
        \SI{57.0}{\milli\litre}. Find its density, and say whether it
        floats on water. \workspace[1.6cm]
  \item Find the mass of \SI{125}{\milli\litre} of ethanol
        ($d = \SI{0.789}{\gram\per\milli\litre}$): \workspace[1.6cm]
  \item A reaction warms a solution by \SI{12}{\celsius}. By how many
        kelvin did it warm? Explain in one line. \workspace[1.4cm]
\end{enumerate}
\selfcheck{(a) \SI{298}{\kelvin}; \SI{77}{\celsius} \quad (b)
\SI{0.789}{\gram\per\milli\litre}, floats \quad (c) \SI{98.6}{\gram}
\quad (d) 12 K --- the degrees are the same size}
\keyonly{\begin{apcanswer}(a) $25 + 273 = \mathbf{298}$~K;
$350 - 273 = \mathbf{77}$~\si{\celsius}.
(b) $d = 45.0/57.0 = \mathbf{0.789}$~g/mL --- less than water's 1.00,
so it \textbf{floats}. (c) $m = dV = 0.789 \times 125 =
\mathbf{98.6}$~g. (d) \textbf{12~K.} A Celsius degree and a kelvin are
the same size, so any temperature \emph{difference} is numerically
identical in the two scales --- only a temperature \emph{value} needs the
273 shifted.\end{apcanswer}}
\end{yourturn}

% =====================================================================
\section*{\sffamily Ladder 7 \textbullet\ Classifying matter}
\zsec{1.10}

\begin{center}
\fbox{\parbox{0.94\linewidth}{\centering
matter $\to$ \textbf{pure substance} (element or compound)
\quad or \quad \textbf{mixture} (homogeneous or heterogeneous)}}
\end{center}

The dividing question is whether the sample has \textbf{fixed
composition}. A compound always has it; a mixture does not --- brass can
be any proportion of copper and zinc.

A \term{homogeneous} mixture (a solution) is uniform throughout;
a \term{heterogeneous} one has visibly distinct regions. And the way you
separate them tells you which you had: mixtures come apart by
\textbf{physical} means (filtration, distillation, chromatography),
compounds only by \textbf{chemical} means.

\begin{workedexample}[I do: classify, then justify by separation]
\textbf{Air} --- a homogeneous \emph{mixture}. Its composition varies
with place and altitude, and it separates by fractional distillation,
which is physical.

\textbf{Water} --- a \emph{compound}. Always 11.2\% hydrogen by mass
wherever it comes from, and splitting it needs electrolysis, which is
chemical.

\textbf{Sand in water} --- a heterogeneous \emph{mixture}: the regions
are visible, and filtration separates them.

\textbf{Copper wire} --- an \emph{element}: it cannot be broken into
anything simpler by any means.

\textbf{The reasoning to imitate:} name the category, then justify it
twice --- once by composition, once by how it separates.
\end{workedexample}

\begin{yourturn}[YOUR TURN 7 --- four questions]
Classify each and give the separation method or reason:
\begin{enumerate}[leftmargin=*,itemsep=0.5em]
  \item table salt, \ce{NaCl} \workspace[1.2cm]
  \item Italian salad dressing \workspace[1.2cm]
  \item filtered seawater \workspace[1.4cm]
  \item Is dissolving sugar in water a physical or chemical change?
        Justify. \workspace[1.4cm]
\end{enumerate}
\selfcheck{(a) compound \quad (b) heterogeneous mixture \quad
(c) homogeneous mixture --- separate by distillation \quad
(d) physical --- the sugar is unchanged and recoverable}
\keyonly{\begin{apcanswer}(a) \textbf{Compound} --- fixed composition,
and only a chemical process (electrolysis of the molten salt) separates
sodium from chlorine. (b) \textbf{Heterogeneous mixture} --- the oil and
vinegar layers are visible, and they separate on standing or with a
separating funnel. (c) \textbf{Homogeneous mixture} --- uniform after
filtering, but the salt content varies by ocean, and simple
\textbf{distillation} recovers the water. (d) \textbf{Physical.} No new
substance forms: the sugar molecules are dispersed but unchanged, and
evaporating the water returns them intact.\end{apcanswer}}
\end{yourturn}

% =====================================================================
\section*{\sffamily Mastery tracker}

Tick a row only if \textbf{all four} YOUR TURN questions were right on
the first attempt.

\renewcommand{\arraystretch}{1.35}
\noindent\begin{tabular}{@{}c p{6.4cm} c p{4.6cm}@{}}
\toprule
\textbf{First try?} & \textbf{Skill} & \textbf{Ladder} &
  \textbf{If not, re-read\ldots}\\
\midrule
$\square$ & units, prefixes, scientific notation & 1 & smaller unit,
  bigger number \\
$\square$ & precision vs.\ accuracy & 2 & the systematic-error case \\
$\square$ & counting significant figures & 3 & the trailing-zero rule \\
$\square$ & calculating with sig figs & 4 & sig figs vs.\ decimal
  places \\
$\square$ & dimensional analysis & 5 & arrange so units cancel \\
$\square$ & temperature and density & 6 & differences vs.\ values \\
$\square$ & classifying matter & 7 & fixed composition? \\
\bottomrule
\end{tabular}

\begin{apcnote}[Scoring yourself honestly]
7/7: you are ready for Chapter 2 --- and, more usefully, for every
calculation in the course. 5--6: redo the missed ladders tomorrow, not
today. 4 or fewer: the misses are almost certainly in Ladders 3--4,
which is worth knowing, because significant figures cost marks in every
later chapter and are the cheapest marks in the course to secure. Redo
those two before anything else.
\end{apcnote}

\end{document}
